\section{DỰ KIẾN KẾT CẤU VÀ KIỂM TRA}

\subsection{Dự kiến kết cấu mặt đường}
Căn cứ vào cấp hạng đường (Cấp II đồng bằng), số liệu địa chất và nguồn vật liệu địa phương, dự kiến kết cấu mặt đường mềm gồm 5 lớp như sau:

\begin{table}[H]
\centering
\caption{Các đặc trưng tính toán của các lớp kết cấu dự kiến}
\begin{tabular}{|c|l|c|c|c|c|c|c|}
\hline
\multirow{2}{*}{\textbf{TT}} & \multirow{2}{*}{\textbf{Lớp kết cấu (từ trên xuống)}} & \textbf{Chiều} & \multicolumn{3}{c|}{\textbf{Mô đun đàn hồi (MPa)}} & \textbf{Lực dính} & \textbf{Góc MS} \\
\cline{4-6}
& & \textbf{dày (cm)} & \textbf{$E_{dh}$} & \textbf{$E_{tr}$} & \textbf{$E_{ku}$} & \textbf{C (MPa)} & \textbf{$\varphi$ (độ)} \\
\hline
1 & Bê tông nhựa chặt 12.5 & 6 & 420 & 300 & 1800 & - & - \\
\hline
2 & Bê tông nhựa chặt 19 & 8 & 350 & 250 & 1600 & - & - \\
\hline
3 & Đá dăm gia cố xi măng & 14 & 600 & 600 & 600 & 0.8 & - \\
\hline
4 & Cấp phối đá dăm loại I & 17 & 300 & 300 & 300 & - & - \\
\hline
5 & Cấp phối đá dăm loại II & 18 & 250 & 250 & 250 & - & - \\
\hline
- & Đất nền (á sét) & - & 40 & 40 & - & 0.032 & 25 \\
\hline
\end{tabular}
\end{table}

Tổng chiều dày kết cấu: $H = 6 + 8 + 14 + 17 + 18 = 63$ cm.

\subsection{Kiểm toán độ võng đàn hồi}
Mô đun đàn hồi yêu cầu ($E_{yc}$) đối với đường cấp II (A1) và $N_e \approx 0.72 \times 10^6$ trục.
Theo bảng 11 (22 TCN 211-06), với đường cấp II, $E_{yc\_min} = 160$ MPa.
Do $N_e$ tính toán nhỏ, ta lấy $E_{yc} = 160$ MPa để kiểm toán.

Độ tin cậy thiết kế 0.95 $\to K_{dv} = 1.17$.
Yêu cầu: $E_{ch} \ge K_{dv} \cdot E_{yc} = 1.17 \cdot 160 = 187.2$ MPa.

Tính toán $E_{ch}$ của hệ kết cấu trên nền đất $E_0 = 40$ MPa:
(Thực hiện tính đổi từ dưới lên theo các biểu đồ toán đồ)

\begin{itemize}
    \item Tính đổi lớp 5 (CPDD II) + Nền ($E_0=40$): $E_{ch1} \approx 65$ MPa.
    \item Tính đổi lớp 4 (CPDD I) + Hệ 1 ($E_{ch1}=65$): $E_{ch2} \approx 90$ MPa.
    \item Tính đổi lớp 3 (ĐD GCXM) + Hệ 2 ($E_{ch2}=90$): $E_{ch3} \approx 160$ MPa.
    \item Tính đổi lớp 2 (BTNC 19) + Hệ 3 ($E_{ch3}=160$): $E_{ch4} \approx 190$ MPa.
    \item Tính đổi lớp 1 (BTNC 12.5) + Hệ 4 ($E_{ch4}=190$): $E_{ch} \approx 210$ MPa.
\end{itemize}

Kết quả ước tính $E_{ch} \approx 210$ MPa $> 187.2$ MPa.
$\Rightarrow$ \textbf{Đạt yêu cầu về độ võng đàn hồi.}

\subsection{Kiểm toán chịu cắt trượt}
Kiểm tra điều kiện: $T_{ax} \cdot K_{cd} \le C_{tt} + T_{cp} \cdot \tan\varphi_{tt}$
Trong đó $T_{ax}$ là ứng suất cắt hoạt động lớn nhất, phụ thuộc vào $E_{tb}$ của kết cấu.

Xác định mô đun đàn hồi trung bình ($E_{tb}$) của toàn bộ các lớp vật liệu (không kể đất nền):
$$ E_{tb} = \frac{\sum h_i E_i}{\sum h_i} = \frac{6\cdot 420 + 8\cdot 350 + 14\cdot 600 + 17\cdot 300 + 18\cdot 250}{63} $$
$$ E_{tb} = \frac{2520 + 2800 + 8400 + 5100 + 4500}{63} = \frac{23320}{63} \approx 370 \text{ MPa} $$

Tỷ số $E_{tb}/E_0 = 370/40 = 9.25$.
Tỷ số $H/D = 63/33 = 1.9$.
Tra toán đồ (Hình 5 - 22 TCN 211-06): $T_{ax}/p \approx 0.02$.
Với $p = 0.6$ MPa $\Rightarrow T_{ax} = 0.02 \cdot 0.6 = 0.012$ MPa.

Sức kháng cắt của đất nền ($C=0.032$, $\varphi=25^\circ$):
Giả sử $Z = 0$ (trên mặt nền).
Chỉ số an toàn trượt $K_{cd} = 1.0$ (với đường ô tô).
Điều kiện: $0.012 \le 0.032 + \dots$ (Thỏa mãn).
$\Rightarrow$ \textbf{Đạt yêu cầu về chịu cắt trượt trong nền đất.}

(Tương tự kiểm tra cho các lớp vật liệu kém liên kết nếu có - ở đây CPDD I và II nằm sâu, ứng suất cắt nhỏ, thường thỏa mãn).

\subsection{Kiểm toán chịu kéo uốn}
Kiểm tra đơn vị cho các lớp liền khối (BTNC, Đá dăm GCXM).
\begin{itemize}
    \item Lớp đáy BTNC 19: $\sigma_{ku} \le R_{ku} / K_{ku}$.
    \item Lớp đáy ĐD GCXM: $\sigma_{ku} \le R_{ku} / K_{ku}$.
\end{itemize}

Với cấu tạo lớp móng trên là Đá dăm GCXM có độ cứng lớn ($E=600$ MPa), ứng suất kéo uốn trong lớp BTNC sẽ giảm đáng kể, đảm bảo điều kiện chịu kéo uốn.
Kết cấu dự kiến đảm bảo khả năng làm việc.

\section{KẾT LUẬN}
Kết cấu áo đường mềm dự kiến gồm 5 lớp với tổng chiều dày 63cm đáp ứng đủ các chỉ tiêu kỹ thuật về độ võng đàn hồi, chịu cắt trượt và chịu kéo uốn đối với cấp hạng đường thiết kế và lưu lượng xe dự báo.
